% use the "wcp" class option for workshop and conference
% proceedings
%\documentclass[gray]{jmlr} % test grayscale version
%\documentclass[tablecaption=bottom]{jmlr}% journal article
\documentclass[pmlr,twocolumn,10pt]{jmlr} % W&CP article

% REQUIRED: select your submission track. Use one of:
  %   proceedings, findings, demo, perspective
  % Leaving the argument empty (\mlhtrack{}) compiles but shows a red
  % "please specify the submission track" warning in the page header.
  % Removing this line entirely raises a "No ML4H track specified" class error.
  % NOTE: keep \mlhtrack{...} BEFORE the \iffinal / \ifmlhneedspmlr blocks below —
  % those conditionals depend on the flags this command sets.
% DECIDED 2026-08-24: Findings (4pp), not Proceedings. Re-checked the
% "auto-cascade" reasoning against the live CFP: decisions for both tracks
% release simultaneously (Oct 22) and both get the same Oct 5-12 rebuttal
% window, so Proceedings carries no actual timeline advantage -- the earlier
% "strictly dominant" framing was wrong on that point. What's left is a real
% fit question: Proceedings wants "high technical sophistication," Findings
% wants "new insights, valuable resources, or exciting preliminary
% directions." A two-test negative result on one task and one architecture
% is a clean Findings paper, not a Proceedings one, and it's unclear whether
% an auto-cascaded submission gets a fresh Findings-appropriate review or
% just carries over Proceedings reviewers' scores against the wrong bar.
% Submitting directly to Findings avoids that risk entirely. NOTE: every
% section below was drafted with an 8pp budget in mind and needs to compress
% by roughly half -- see the per-section space notes added at fix-in-place
% time. "Lessons for the Field" is now a CANDIDATE FOR MERGING into
% Section~\ref{sec:confounds} at prose time rather than standing alone.
\mlhtrack{findings}

\newif\iffinal
\finalfalse  % toggle this flag for initial submission
% \finaltrue  % toggle this flag for camera-ready version after acceptance

\iffinal
    % For Proceedings and Perspective tracks (published in PMLR), replace XXX below
    % with the specific PMLR volume number sent to you before the camera-ready submission.
    \ifmlhneedspmlr
      \jmlrvolume{XXX}
      \jmlryear{2026}
    \fi
    \ifmlhfindings \jmlrproceedings{}{ML4H 2026 - Findings Track}\fi
    \ifmlhdemo     \jmlrproceedings{}{ML4H 2026 - Demo Track}\fi
    \jmlrworkshop{Machine Learning for Health (ML4H) 2026}
\else
    % Submitted for review
    \jmlrproceedings{}{Submitted to ML4H 2026: \mlhtrackname}
    \jmlrworkshop{Machine Learning for Health (ML4H) 2026}
\fi

% \usepackage{geometry}
% \geometry{margins=0.1in,textwidth=7in}

% The following packages will be automatically loaded:
% amsmath, amssymb, natbib, graphicx, url, algorithm2e
\urlstyle{same}

%\usepackage{rotating}% for sideways figures and tables
%\usepackage{longtable}% for long tables

% The booktabs package is used by this sample document
% (it provides \toprule, \midrule and \bottomrule).
% Remove the next line if you don't require it.

\usepackage{booktabs}
% The siunitx package is used by this sample document
% to align numbers in a column by their decimal point.
% Remove the next line if you don't require it.
\usepackage{siunitx}

% The lineno package is required for denoting line
% numbers for paper review.
\usepackage[switch]{lineno}

% The following command is just for this sample document:
\newcommand{\cs}[1]{\texttt{\char`\\#1}}% remove this in your real article

% The following is to recognise equal contribution for authorship
\newcommand{\equal}[1]{{\hypersetup{linkcolor=black}\thanks{#1}}}

% Define an unnumbered theorem just for this sample document for
% illustrative purposes:
\theorembodyfont{\upshape}
\theoremheaderfont{\scshape}
\theorempostheader{:}
\theoremsep{\newline}
\newtheorem*{note}{Note}

% The optional argument of \title is used in the header
% Working title -- open to a better one. Short-title arg is what appears in
% the running header / TOC, must stay under ~50 chars.
\title[Why a Principled Physiology Constraint Failed]{Plausible, Scrutinized, Dismantled: A Rigorous Negative Result for Physiology-Constrained Cross-Hospital Transfer}

% Anything in the title that should appear in the main title but 
% not in the article's header or the volume's table of
% contents should be placed inside \titletag{}

%\title{Title of the Article\titletag{\thanks{Some footnote}}}


% Use \Name{Author Name} to specify the name.
% If the surname contains spaces, enclose the surname
% in braces, e.g. \Name{John {Smith Jones}} similarly
% if the name has a "von" part, e.g \Name{Jane {de Winter}}.
% If the first letter in the forenames is a diacritic
% enclose the diacritic in braces, e.g. \Name{{\'E}louise Smith}

% \thanks must come after \Name{...} not inside the argument for
% example \Name{John Smith}\nametag{\thanks{A note}} NOT \Name{John
% Smith\thanks{A note}}

% Anything in the name that should appear in the title but not in the 
% article's header or footer or in the volume's
% table of contents should be placed inside \nametag{}

% Anonymized author list during review (ML4H is double-blind).
% *** BLOCKING ITEM, 2026-08-24: Rahmani's co-authorship is NOT confirmed. ***
% Every prior interaction with him on either PCL paper has been mentor/advisor,
% never co-author. His name is held here as a placeholder ONLY -- do not
% treat it as settled, do not move to full prose past the Introduction's
% framing sections until he has given explicit yes/no on co-authorship.
% CURRENT WORKING ASSUMPTION (pending his answer): NOT onboard. Fallback if
% he declines or doesn't respond in time: solo byline (Nikhil Tamvada,
% Independent Researcher) with Rahmani thanked in \acks{} instead.
% De-anonymize only in the camera-ready version if accepted; keep
% "Anonymous Author(s)" for every review-stage submission regardless of how
% this resolves.
\author{Anonymous Author(s)}

% Two authors with the same address
% \author{%
%  \Name{Author Name1\nametag{\thanks{A note}}} \Email{abc@sample.com}\and
%  \Name{Author Name2} \Email{xyz@sample.com}\\
%  \addr Address
% }

% Three or more authors with the same address:
% \author{%
%  \Name{Author Name1} \Email{an1@sample.com}\\
%  \Name{Author Name2} \Email{an2@sample.com}\\
%  \Name{Author Name3} \Email{an3@sample.com}\\
%  \Name{Author Name4} \Email{an4@sample.com}\\
%  \Name{Author Name5} \Email{an5@sample.com}\\
%  \Name{Author Name6} \Email{an6@sample.com}\\
%  \Name{Author Name7} \Email{an7@sample.com}\\
%  \Name{Author Name8} \Email{an8@sample.com}\\
%  \Name{Author Name9} \Email{an9@sample.com}\\
%  \Name{Author Name10} \Email{an10@sample.com}\\
%  \Name{Author Name11} \Email{an11@sample.com}\\
%  \Name{Author Name12} \Email{an12@sample.com}\\
%  \Name{Author Name13} \Email{an13@sample.com}\\
%  \Name{Author Name14} \Email{an14@sample.com}\\
%  \addr Address
% }

% % Authors with different addresses and equal first authors:
% \author{%
% \Name{First Author 1}\equal{These authors contributed equally} \Email{abc@sample.com}\\
% \addr University X, Country 1
% \AND
% % footnotemark[1] is to refer to the \equal footnote
% \Name{First Author 2}\footnotemark[1] \Email{def@sample.com}\\
% \addr University Y, Country 2
% \AND
% \Name{Last Author} \Email{ghi@sample.com}\\
% \addr University Z, Country 3
% }

%%%%%%%%%%%%%%%%%%%%%%%%%%%%%%%%%%%%%%%%%%%%%%%%%%%%%%%%%%%%%%%%%%%%%%%%
%%%%%%%%%%%%% Remove the \linenumbers in the final version %%%%%%%%%%%%%
%%%%%%%%%%%%%%%%%%%%%%%%%%%%%%%%%%%%%%%%%%%%%%%%%%%%%%%%%%%%%%%%%%%%%%%%
\linenumbers % Activate line numbering

\begin{document}

\maketitle

%%%%%%%%%%%%%%%%%%%%%%%%%%%%%%%%%%%%%%%%%%%%%%%%%%%%%%%%%%%%%%%%%%%%%%%%
%%%%% FULL PROSE, 2026-08-24. Numbers sourced from                  %%%%%
%%%%% _archive/pcl_original_draft/fmain_type2.tex (original result)  %%%%%
%%%%% and chat2_papers/paper/build_urtc.py (corrected result, all    %%%%%
%%%%% five-confound text, IRM/XGBoost provenance). Byline still      %%%%%
%%%%% pending Rahmani's answer -- see \author comment above; nothing %%%%%
%%%%% below depends on how that resolves except \acks.              %%%%%
%%%%%%%%%%%%%%%%%%%%%%%%%%%%%%%%%%%%%%%%%%%%%%%%%%%%%%%%%%%%%%%%%%%%%%%%

\ifmlhdemo\else
\begin{abstract}
Physiology-constrained pretraining does not improve cross-hospital sepsis
prediction. A well-motivated hypothesis -- that penalizing violations of
known physiological relationships during masked-value pretraining of a
clinical time-series transformer would yield representations that transfer
better across hospitals than an unconstrained baseline -- initially appeared
confirmed: the constrained model beat empirical risk minimization (ERM) at
every held-out site. Systematic re-evaluation traced the entire apparent
gain to five identifiable evaluation confounds. Correcting all five reverses
the result: the constrained model falls below ERM at every site, and a plain
gradient-boosted-tree baseline on aggregate features matches or beats both
neural methods. Run-to-run seed variance of 4 to 7 AUROC points exceeds most
of the originally reported effect sizes. The
negative result also holds under a second, independently designed test: a
label-free model-selection criterion built on the same physiological signal
fails to reliably outperform a simple unsupervised baseline. We report each
confound mechanistically and in generalizable terms, so a reader building a
different domain-knowledge-constrained clinical method can recognize the
same traps before publishing them.
\end{abstract}
\begin{keywords}
physiology-constrained learning, negative results, cross-hospital
generalization, evaluation confounds, clinical machine learning,
reproducibility
\end{keywords}
\fi

\ifmlhneedsstatements
\paragraph*{Data and Code Availability}
All experiments use three publicly available, de-identified,
credentialed-access ICU databases: the PhysioNet/Computing in Cardiology
Challenge 2019 dataset \citep{reyna2020early}, MIMIC-IV
\citep{johnson2023mimiciv}, and the eICU Collaborative Research Database
\citep{pollard2018eicu}. Code -- the corrected SOFA-based label-derivation
and source-domain-only evaluation pipeline central to
Section~\ref{sec:confounds} -- will be released as an anonymized
repository at submission.
% PRE-SUBMISSION TODO, 2026: this commitment is not yet fulfilled -- an
% anonymized repo/zip actually needs to exist before submission, per the
% template's own rule that unanonymized or missing code at submission time
% is an automatic desk-reject. Needs: the SOFA label-derivation code
% (chat1_protocol/src/data/sofa_sepsis.py) and the corrected eval pipeline
% (chat2_papers/paper/build_urtc.py's underlying training/eval code, not
% the docx-builder script itself), stripped of author-identifying paths
% and comments, in a standalone anonymized location.

\paragraph*{Institutional Review Board (IRB)}
This work is a secondary analysis of pre-existing, de-identified,
credentialed-access datasets already approved for secondary research use
under their respective data use agreements. No new human-subjects data was
collected, and no additional IRB approval was required.
\fi

\section{Introduction}
\label{sec:intro}

Clinical machine learning models trained at one hospital routinely lose
accuracy when deployed at another, and the consequences are not
hypothetical. An external validation of a proprietary sepsis prediction
model widely deployed across US hospitals found an AUROC of 0.63, a
sensitivity of 33\%, and a positive predictive value of 12\% at a single
academic medical center -- well below the vendor's reported performance --
contributing to the model's deactivation at that site
\citep{wong2021external}.

This paper tests a specific, principled response: physiology-constrained
learning (PCL), which augments masked-value pretraining of a clinical
time-series transformer with differentiable penalties for violating known
physiological relationships among simultaneously measured variables. The
premise follows the Independent Causal Mechanisms principle
\citep{scholkopf2021toward}: relationships such as the mean arterial
pressure identity or Henderson-Hasselbalch hold at every hospital
regardless of local measurement device, calibration, or ordering practice,
so a representation constrained to respect them should depend less on
site-specific artifacts. This is not a strawman; it is the response a
careful practitioner reaches for when building a model that generalizes
across sites, and the hypothesis an earlier version of this investigation
set out to confirm.

Our contribution is not the claim that this hypothesis failed; it is a
rigorous diagnostic accounting of why a principled hypothesis produced a
false positive convincing enough to submit, general enough that a reader
building a different domain-knowledge-constrained method can recognize the
same traps before publication. An earlier version of this investigation was
submitted elsewhere and withdrawn after our own follow-up uncovered the
confounds this paper documents; we report that investigation in full here.

\section{Method}
\label{sec:method}

PCL pretrains a six-layer clinical time-series transformer ($d{=}256$, 8
heads, $\sim$3.2M parameters) on 48-hour, hourly-binned ICU windows using a
masked-value prediction objective, following the general masked-pretraining
paradigm \citep{vaswani2017attention,devlin2019bert}. The training objective
is $L = L_{\mathrm{mask}} + \lambda L_{\mathrm{phys}}$, where
$L_{\mathrm{phys}}$ averages differentiable residuals on three physiological
relations -- an $L_1$ penalty on the mean arterial pressure identity and
squared-error penalties on the Henderson-Hasselbalch and Severinghaus
equations -- applied only at timesteps where all required variables were
observed and at least one was masked, so the penalty must be satisfied by
imputation from cross-variable structure rather than by copying the input.
Models are trained on PhysioNet Site A (10,381 stays, 3.3\% sepsis
prevalence) and evaluated zero-shot on PhysioNet Site B (14,779 stays,
2.1\%), MIMIC-IV (74,607 stays, 30.9\%), and eICU-CRD (130,446 stays,
8.8\%) on the task of sepsis onset prediction. The corrected pipeline
(Section~\ref{sec:corrected}) uses 17 physiological variables, up from the
9 used in the original result reported next -- the eight added (temperature,
respiratory rate, white blood cell count, lactate, BUN, creatinine,
potassium, glucose) are standard ICU vitals and routine
chemistry/hematology, adopted in the same commit as the corrected SOFA
labeling, not tuned after seeing corrected-pipeline results, so this
expansion is not itself a sixth confound. It improved external transfer for
both ERM and PCL and was retained throughout, so Table~\ref{tab:initial}
versus Table~\ref{tab:corrected} compares two different input
representations, not the same model re-scored under corrected labels
alone.

\section{The Initial, Misleading Result}
\label{sec:initial}

Table~\ref{tab:initial} reports the original, first-pass result: PCL appears
to beat every baseline at every held-out site, with the margin growing with
distributional distance -- +6.4 percentage points at Site B, +10.1 at
MIMIC-IV, +8.4 at eICU -- and PCL is the best method in every one of three
seeds at every site. The result also passed a constraint-randomization
ablation (real physiological constraints outperformed shuffled,
wrong-equation, and noise-injected variants, with the two most-corrupted
variants falling below unconstrained ERM) and a measurement-noise robustness
check (PCL's AUROC degraded less than ERM's under injected sensor bias).
Nothing about the result looked fragile from any angle tried at the time;
that is precisely what makes the diagnosis in Section~\ref{sec:confounds}
worth reporting in full.

\begin{table}[htbp]
\floatconts
  {tab:initial}%
  {\caption{The original, uncorrected result: zero-shot cross-hospital sepsis
   AUROC, mean over 3 seeds. PCL appears to beat every baseline at every
   held-out site. This result does not survive Section~\ref{sec:confounds}
   and is reported here only to show how convincing it looked.}}%
  {\begin{tabular}{lcccc}
  \toprule
  Method & Source & Site B & MIMIC-IV & eICU \\
  \midrule
  ERM & 0.823 & 0.697 & 0.582 & 0.541 \\
  DRO & 0.778 & 0.683 & 0.612 & 0.604 \\
  IRM$^\dagger$ & --- & 0.675 & --- & --- \\
  \textbf{PCL} & \textbf{0.865} & \textbf{0.761} & \textbf{0.682} & \textbf{0.626} \\
  \bottomrule
  \end{tabular}}
\end{table}
[$^\dagger$IRM: single seed, Site B only, $k{=}3$ care-unit environments.
Margins as originally reported: +6.4pp Site B, +10.1pp MIMIC-IV, +8.4pp eICU.]

\section{The Five Confounds}
\label{sec:confounds}

Each confound below is stated mechanistically and in terms general enough
that a reader building a different domain-knowledge-constrained method can
check their own pipeline against it; each ends with the generalized lesson.

\subsection{Label-Definition Shift}
Training labels followed the clinical Sepsis-3 consensus definition
\citep{singer2016sepsis3}, while external evaluation on MIMIC-IV and eICU
originally used administrative ICD-code-derived sepsis labels. Clinical and
claims-based sepsis definitions are known to diverge substantially
\citep{rhee2017incidence}, so cross-site AUROC differences conflated genuine
distribution shift with a change in the definition of the outcome itself.
We corrected this by re-deriving Sepsis-3 directly via SOFA scoring anchored
to a suspected-infection window, following the standard reference
implementation \citep{seymour2016assessment}, yielding sepsis prevalence of
30.9\% on MIMIC-IV and 8.8\% on eICU. \textbf{Lesson:} any cross-site
comparison must confirm the outcome definition, not only the input
distribution, is held fixed across sites.

\subsection{Pretraining Data Leakage Disguised as Zero-Shot Transfer}
PCL's masked-value pretraining used unlabeled data from all four sites,
while ERM's pretraining used source-site data only. Although pretraining is
self-supervised and touches no labels, exposure to target-site inputs during
pretraining still breaks the zero-shot premise the comparison depends on:
PCL had, in effect, already seen the deployment distribution. The corrected
pipeline pretrains PCL source-only, matching ERM. \textbf{Lesson:} ``no labels touched'' is not the same guarantee as ``no
target data touched''; a zero-shot claim must audit every stage of
training, not only the labeled fine-tuning stage.

\subsection{OOD-Contaminated Hyperparameter Selection}
The constraint weight $\lambda$, which strongly affects cross-hospital
performance, was originally chosen using AUROC on a held-out target site
rather than on source-domain validation data, so the reported zero-shot
numbers were selected with knowledge of the test distribution they were
meant to predict. \textbf{Lesson:} a hyperparameter is only as zero-shot as
the data used to select it; auditing final-model performance is not
sufficient if the selection procedure that produced that model touched
target data.

\subsection{Partial Circularity in the Severinghaus Constraint}
PhysioNet does not record arterial oxygen partial pressure (PaO2) directly;
the released dataset reconstructs it from oxygen saturation by inverting the
Severinghaus equation -- the same relation PCL's constraint penalizes the
model for violating. The constraint's training-time supervision was
therefore partly self-referential: it could be satisfied, in part, by
recovering a value the data-generating process had already computed from the
constraint's own equation. \textbf{Lesson:} before enforcing a
domain-knowledge constraint, audit whether any of its input variables were
themselves derived by assuming that constraint; a preprocessing pipeline can
silently manufacture the invariance a model is later credited with
discovering.

\subsection{Measurement-Composition Artifact in the Aggregate Violation Score}
The aggregate physiological violation score averages whichever of the three
constraints happen to be computable at a given timestep, and constraint
availability is highly uneven and site-dependent -- for instance,
bicarbonate is charted in 8.1\% of Site A stay-hours versus 0.16\% at Site
B, a fifty-fold difference driven by local lab-ordering practice, not
physiology. Because the sparser constraints also occupy a different residual
scale than the dense ones, a site's aggregate mean partly tracks which labs
are ordered there rather than any physiological difference. Recomputing
each site's aggregate with pooled component values (raw data, no model)
isolates roughly 49\% of the gap as composition artifact and 51\% as
genuine -- partially, not wholly, artifactual. \textbf{Lesson:} an aggregate score
over heterogeneously-available components is a composition statistic
before it is anything else; decompose it before reading a cross-site
difference in it as a physiological effect.

\section{The Corrected Result}
\label{sec:corrected}

With all five confounds addressed, PCL falls below ERM at every site
(Table~\ref{tab:corrected}), alongside a group-distributionally-robust
baseline \citep{sagawa2020distributionally}. The Source column also moves
for both methods, not only the targets; expected, since the 9-to-17
expansion (Section~\ref{sec:method}) and the new $\lambda$ procedure
(below) make the corrected PCL a genuinely different trained model, not
the same one re-scored.

A gradient-boosted tree \citep{chen2016xgboost} trained on simple per-stay
summary statistics matches or exceeds both neural methods at every external
site (0.754 / 0.654 / 0.660 versus ERM's 0.701 / 0.649 / 0.592 and PCL's
0.655 / 0.622 / 0.572), which further questions whether the transformer
architecture is earning its complexity here; it has no Source entry because
it is used only as an external baseline, never a source-selection
candidate. Table~\ref{tab:corrected} omits IRM entirely -- unlike its
$k{=}3$ single-seed appearance in Table~\ref{tab:initial}, no version of it
was ever run zero-shot cross-hospital. We separately checked whether
invariant risk minimization \citep{arjovsky2019invariant} with 30 real
hospital environments drawn from eICU closes the gap to ERM; it does not,
but that check turned out not to support any claim about this paper's zero-shot
setting, for two reasons, and we report it only for completeness. First,
that comparison is in-distribution within eICU, not zero-shot
Site-A-to-eICU transfer, so it does not speak to the transfer failure this
paper is about. Second, the gap it found (0.788 versus its own ERM baseline
of 0.792, 0.4 percentage points) is smaller than this study's own
run-to-run seed noise (4 to 7 percentage points, below), so it is not
distinguishable from noise in the first place. We do not treat this as
evidence either way.

Run-to-run variability across seeds reached 4 to 7 AUROC points on this
pipeline, exceeding most of the differences originally reported as method
effects in Table~\ref{tab:initial}: a single-seed comparison at this scale
is not informative regardless of how large the claimed gain looks.

The constraint weight $\lambda$ itself was selected by sweeping six
candidate values and taking the one with the highest AUROC on a held-out
validation split of the source site (PhysioNet Site A) only -- the same
split used for every other hyperparameter decision, never a target site.
This replaces the original selection procedure named in
Section~\ref{sec:confounds}.

\begin{table}[htbp]
\floatconts
  {tab:corrected}%
  {\caption{Corrected cross-hospital sepsis AUROC, mean over 3 seeds, all
   five confounds addressed. PCL is now below ERM at every site.}}%
  {\begin{tabular}{lcccc}
  \toprule
  Method & Source & Site B & MIMIC-IV & eICU \\
  \midrule
  ERM & 0.818 & 0.701 & 0.649 & 0.592 \\
  PCL & 0.798 & 0.655 & 0.622 & 0.572 \\
  Group DRO & 0.822 & 0.717 & 0.688 & 0.628 \\
  XGBoost & --- & 0.754 & 0.654 & 0.660 \\
  \bottomrule
  \end{tabular}}
\end{table}

\section{Beyond Prediction Accuracy}
\label{sec:beyond}

The same physiological signal was also tested as a label-free criterion for
selecting the constraint weight without target labels -- an independently
designed test, since a selection criterion can fail even when the
underlying representation does not. Prior work by the same author
\citep{tamvada2026urtc} finds it does not reliably outperform a simple
unsupervised baseline (reconstruction error). That work is itself
unpublished and under review, so we report only the qualitative outcome,
not its numbers, and do not reproduce its tables; see it for the full
comparison.

A calibration or conformal-prediction-coverage test was considered but not
included: no such artifact was recoverable, and we chose not to report an
unverifiable number.

\section{Limitations}
\label{sec:limitations}

This is a case study in how one principled domain-knowledge constraint
failed on one task and one architecture, not a universal claim that such
constraints cannot improve cross-hospital transfer. All experiments use a
single task (sepsis onset) and a single architecture and pretraining
objective; whether the same confounds recur elsewhere is untested. Three
seeds is a small sample for a variance estimate, though the seed-variance
finding itself (Section~\ref{sec:corrected}) is the point, not a caveat.
We also did not decompose each confound's individual contribution to the
Table~\ref{tab:initial}-to-\ref{tab:corrected} gap: isolating one needs
retraining with it alone reintroduced, and at this study's 4-to-7-point
seed-noise floor a single confound's effect may not be separable from
noise, so we report the five jointly rather than rank them.

% NOTE 2026-08-24: if Rahmani declines/doesn't respond, add a second
% acknowledgment sentence here thanking him for advisory input (fallback
% plan noted at the \author block above). Currently dataset-thanks only.
%
% FIX 2026-08-24, found by compiling: the template's own instruction says
% acknowledgments "should only appear in the camera-ready version of the
% paper if it is accepted," but \acks (jmlr.cls:595, just
% \section*{Acknowledgments}#1) has no internal gating on the \iffinal
% toggle already defined in the preamble -- so left unconditional, it
% renders during review too, eating page-limit-relevant space it isn't
% supposed to use. Gating it here is also what fixed a real Acknowledgments/
% References layout collision on page 4 of the compiled proxy: with \acks
% removed from the review-stage page count, everything before References
% now has room to breathe. Toggle \finaltrue (top of file) to see it.
\iffinal
\acks{We thank the maintainers of PhysioNet, MIMIC-IV, and eICU-CRD for
making these datasets available for secondary research use.}
\fi

\bibliography{ref}

% deadpcl/ref.bib still needs to be created (does not exist yet -- \bibliography
% will fail to resolve citations until it does). Every \citep/\citet key used
% above, all independently verified (see chat transcript 2026-08-24 for
% source of each), needed as a bibitem:
%   wong2021external        - Wong et al., JAMA Intern Med 181(8):1065-1070, 2021
%   scholkopf2021toward      - Scholkopf et al., Proc. IEEE 109(5):612-634, 2021
%   singer2016sepsis3        - Singer et al., JAMA 315(8):801-810, 2016
%   seymour2016assessment    - Seymour et al., JAMA 315(8):762-774, 2016
%   rhee2017incidence        - Rhee et al., JAMA 318(13):1241-1249, 2017
%   reyna2020early           - Reyna et al., Crit. Care Med. 48(2):210-217, 2020
%   johnson2023mimiciv       - Johnson et al., Sci. Data 10(1), 2023
%   pollard2018eicu          - Pollard et al., Sci. Data 5, 2018
%   vaswani2017attention     - Vaswani et al., NeurIPS 2017
%   devlin2019bert           - Devlin et al., NAACL-HLT 2019, pp. 4171-4186
%   chen2016xgboost          - Chen & Guestrin, ACM SIGKDD 2016, pp. 785-794
%   arjovsky2019invariant    - Arjovsky et al., arXiv:1907.02893, 2019
%   sagawa2020distributionally - Sagawa et al., ICLR 2020
%   rosenfeld2021risks       - Rosenfeld et al., ICLR 2021
%   tamvada2026urtc          - SELF-CITATION, this author's URTC paper,
%                              "When Validation Lies: Diagnosing and Fixing
%                              Model Selection Failures Under Real Hospital
%                              Shift," submitted -- exact venue/year string
%                              for the bibitem needs confirming once its own
%                              submission status is final.
%
% The confound-detection paper is deliberately NOT a \citep in this draft --
% it has no formal title yet (chat1_protocol/README.md describes it only as
% "a new methods paper"), so it's referenced in prose only
% (Section~\ref{sec:confounds}'s closing sentence) rather than invented as a
% bibitem. Add a real citation only once that paper has an actual title.

\end{document}
